
\documentclass[tikz,border=8pt]{standalone}
\usepackage[UTF8,fontset=fandol]{ctex}
\usepackage{amsmath,amssymb}
\usetikzlibrary{arrows.meta,calc,positioning}
\definecolor{ink}{HTML}{172033}
\definecolor{muted}{HTML}{637083}
\definecolor{blue}{HTML}{4F86C6}
\definecolor{green}{HTML}{61A66A}
\definecolor{red}{HTML}{D55E5E}
\definecolor{amber}{HTML}{D59B3D}
\definecolor{paperblue}{HTML}{F2F7FD}
\definecolor{papergreen}{HTML}{F3F9F1}
\definecolor{paperamber}{HTML}{FFF8E8}
\tikzset{
  panel/.style={anchor=north west,minimum width=9.15cm,minimum height=6.05cm,
    text width=8.45cm,align=left,inner sep=10pt,draw=black!12,rounded corners=5pt,fill=white},
  title/.style={font=\bfseries\large,text=ink,anchor=west},
  body/.style={font=\small,text=muted,align=left,text width=7.9cm},
  formula/.style={draw=blue!35,fill=paperblue,rounded corners=4pt,inner sep=8pt,
    text=ink,align=center,text width=7.7cm},
  insight/.style={draw=green!40,fill=papergreen,rounded corners=4pt,inner sep=7pt,
    text=ink,align=left,text width=7.7cm,font=\small},
  arrow/.style={-{Latex[length=2.3mm]},line width=1pt,draw=blue!70},
  chip/.style={rounded corners=3pt,draw=black!14,fill=black!2,inner xsep=7pt,
    inner ysep=5pt,font=\small,text=ink},
}
\newcommand{\Panel}[5]{%
  \node[panel] (#1) at (#2,#3) {};
  \node[title] at ($(#1.north west)+(0.35,-0.42)$) {#4};
  \node[body,anchor=north west] at ($(#1.north west)+(0.35,-1.05)$) {#5};
}
\newcommand{\Summary}[1]{%
  \node[draw=amber!35,fill=paperamber,rounded corners=6pt,minimum width=28.2cm,
    minimum height=1.55cm,text width=27.2cm,align=center,font=\small,text=ink]
    at (14.125,-13.35) {\textbf{一图总结}\quad #1};
}
\newcommand{\Identity}{%
  \node[font=\scriptsize,text=black!38,anchor=east] at (28.15,-14.35)
    {cs229-storyboard@五道口纳什};
}
\newcommand{\Formula}[1]{%
  \par\vspace{5pt}\begingroup\setlength{\fboxsep}{8pt}%
  \colorbox{paperblue}{\parbox{7.45cm}{\centering\color{ink}#1}}%
  \endgroup\par\vspace{6pt}}
\newcommand{\Insight}[1]{%
  \par\vspace{5pt}\begingroup\setlength{\fboxsep}{7pt}%
  \colorbox{papergreen}{\parbox{7.55cm}{\color{ink}#1}}%
  \endgroup\par\vspace{6pt}}
\newcommand{\Chip}[1]{%
  \begingroup\setlength{\fboxsep}{5pt}\colorbox{black!4}{\strut\color{ink}#1}\endgroup}

\begin{document}
\begin{tikzpicture}[x=1cm,y=1cm]
\fill[white] (-0.35,0.35) rectangle (28.6,-14.55);
\Panel{p1}{0}{0}{1. K-means 先分配}{固定中心，选择最近簇：\Formula{$c_i=\arg\min_k\lVert x_i-\mu_k\rVert_2^2$}}
\Panel{p2}{9.55}{0}{2. 再更新中心}{固定分配，取簇内均值：\Formula{$\mu_k=\dfrac{\sum_i\mathbf1[c_i=k]x_i}{\sum_i\mathbf1[c_i=k]}$}两步交替降低平方距离目标。}
\Panel{p3}{19.1}{0}{3. 硬分配有明显限制}{\Insight{K-means 偏好近似球形、相似尺度的簇，并且不给出不确定性。初始中心不同可能得到不同结果。}}
\Panel{p4}{0}{-6.45}{4. GMM 把簇变成概率分布}{\Formula{$p(x)=\sum_{k=1}^K\pi_k\mathcal N(x;\mu_k,\Sigma_k)$}每个样本可以部分属于多个簇。}
\Panel{p5}{9.55}{-6.45}{5. E 步计算软责任度}{\Formula{$\gamma_{ik}=\dfrac{\pi_k\mathcal N(x_i;\mu_k,\Sigma_k)}{\sum_j\pi_j\mathcal N(x_i;\mu_j,\Sigma_j)}$}}
\Panel{p6}{19.1}{-6.45}{6. M 步做加权估计}{\Formula{$N_k=\sum_i\gamma_{ik},\quad\mu_k=\dfrac1{N_k}\sum_i\gamma_{ik}x_i$}协方差和混合权重同样按责任度更新。}
\Summary{K-means 硬分配 $\leftrightarrow$ 中心更新 $\rightarrow$ GMM 概率建模 $\rightarrow$ E 步责任度 $\leftrightarrow$ M 步参数更新。}
\Identity
\end{tikzpicture}
\end{document}
